% Options for packages loaded elsewhere
\PassOptionsToPackage{unicode}{hyperref}
\PassOptionsToPackage{hyphens}{url}
\PassOptionsToPackage{dvipsnames,svgnames,x11names}{xcolor}
%
\documentclass[
  10pt,
]{article}
\usepackage{amsmath,amssymb}
\usepackage{iftex}
\ifPDFTeX
  \usepackage[T1]{fontenc}
  \usepackage[utf8]{inputenc}
  \usepackage{textcomp} % provide euro and other symbols
\else % if luatex or xetex
  \usepackage{unicode-math} % this also loads fontspec
  \defaultfontfeatures{Scale=MatchLowercase}
  \defaultfontfeatures[\rmfamily]{Ligatures=TeX,Scale=1}
\fi
\usepackage{lmodern}
\ifPDFTeX\else
  % xetex/luatex font selection
  \setmainfont[]{DejaVu Serif}
  \setsansfont[]{DejaVu Sans}
  \setmonofont[]{DejaVu Sans Mono}
  \setmathfont[]{Latin Modern Math}
\fi
% Use upquote if available, for straight quotes in verbatim environments
\IfFileExists{upquote.sty}{\usepackage{upquote}}{}
\IfFileExists{microtype.sty}{% use microtype if available
  \usepackage[]{microtype}
  \UseMicrotypeSet[protrusion]{basicmath} % disable protrusion for tt fonts
}{}
\makeatletter
\@ifundefined{KOMAClassName}{% if non-KOMA class
  \IfFileExists{parskip.sty}{%
    \usepackage{parskip}
  }{% else
    \setlength{\parindent}{0pt}
    \setlength{\parskip}{6pt plus 2pt minus 1pt}}
}{% if KOMA class
  \KOMAoptions{parskip=half}}
\makeatother
\usepackage{xcolor}
\usepackage[margin=2.0cm]{geometry}
\usepackage{longtable,booktabs,array}
\usepackage{calc} % for calculating minipage widths
% Correct order of tables after \paragraph or \subparagraph
\usepackage{etoolbox}
\makeatletter
\patchcmd\longtable{\par}{\if@noskipsec\mbox{}\fi\par}{}{}
\makeatother
% Allow footnotes in longtable head/foot
\IfFileExists{footnotehyper.sty}{\usepackage{footnotehyper}}{\usepackage{footnote}}
\makesavenoteenv{longtable}
\setlength{\emergencystretch}{3em} % prevent overfull lines
\providecommand{\tightlist}{%
  \setlength{\itemsep}{0pt}\setlength{\parskip}{0pt}}
\setcounter{secnumdepth}{-\maxdimen} % remove section numbering
% doc_header.tex -- LaTeX preamble additions used by build_docs.sh for every
% Markdown -> LaTeX -> PDF conversion (supplements, README, submission files).
\usepackage{etoolbox}
% longtable may not be loaded by pandoc when a document has no table
\usepackage{longtable}
% bold vectors: \bm -> \symbf under unicode-math (XeLaTeX), bm package otherwise
\makeatletter
\@ifpackageloaded{unicode-math}{\providecommand{\bm}[1]{\symbf{#1}}}{\usepackage{bm}}
\makeatother
% smaller type in long tables so that wide audit tables fit the text width;
% inside a table a line break is allowed after an underscore of plain text
% (pandoc writes it as \_), so that long identifiers wrap inside their column
% rows are set a little apart, so that symbols with indices do not touch
\AtBeginEnvironment{longtable}{\footnotesize\renewcommand{\arraystretch}{1.45}%
  \let\sptunderscore\_\def\_{\sptunderscore\allowbreak}}
\setlength{\LTcapwidth}{\textwidth}
% allow line breaks after slashes and hyphens inside table cells
\usepackage{xurl}
% code blocks: break over-long lines instead of running into the margin.
% Plain ``` blocks are typeset by pandoc as {verbatim}, highlighted ones
% (```bash) as {Highlighting}; both are re-declared as breakable fancyvrb
% environments.  The continuation line is marked with a small hooked arrow.
\usepackage{fvextra}
\RecustomVerbatimEnvironment{verbatim}{Verbatim}{breaklines,breakanywhere,fontsize=\small}
\makeatletter
\@ifundefined{Highlighting}{}{%
  \RecustomVerbatimEnvironment{Highlighting}{Verbatim}{breaklines,breakanywhere,fontsize=\small,commandchars=\\\{\}}}
\makeatother
% inline code (pandoc writes it as \texttt{...}): never hyphenate -- an inserted
% hyphen could be read as part of a file name -- but allow a line break after
% an underscore, slash, comma, colon or equals sign, so that long paths and
% column lists wrap inside the text width and inside table columns.
% Inside a formula (e.g. the name of a notebook symbol in an equation) the same
% command must first leave math mode, which \text of amsmath does.
\ExplSyntaxOn
\tl_new:N \l__spt_code_tl
\cs_new_protected:Npn \__spt_texttt:n #1
  {
    \tl_set:Nn \l__spt_code_tl {#1}
    \regex_replace_all:nnN { ([\/\,\:\=]) } { \1 \c{allowbreak} } \l__spt_code_tl
    \regex_replace_all:nnN { \c{_} } { \c{_} \c{allowbreak} } \l__spt_code_tl
    \group_begin: \ttfamily \hyphenchar \font = -1 \scan_stop: \l__spt_code_tl \group_end:
  }
\RenewDocumentCommand \texttt { m }
  {
    \mode_if_math:TF { \text { \__spt_texttt:n {#1} } } { \__spt_texttt:n {#1} }
  }
\ExplSyntaxOff
% A figure is written in the Markdown source as an image paragraph followed by
% a caption paragraph ("**Figure S1.** ..."), so that the caption is also
% visible on GitHub.  An image that starts a paragraph is centered.
\usepackage{graphicx}% pandoc loads it only for documents that contain an image
\NewCommandCopy\sptincludegraphics\includegraphics
\RenewDocumentCommand\includegraphics{O{} m}{%
  \ifvmode\begin{center}\sptincludegraphics[#1]{#2}\end{center}%
  \else\sptincludegraphics[#1]{#2}\fi}
% keep figures where they are referenced
\usepackage{float}
\floatplacement{figure}{H}
% running header with the document title
\usepackage{fancyhdr}
\pagestyle{fancy}
\fancyhf{}
\fancyhead[L]{\small PROJECT spacetimes, Revision\_Deflation\_CJT\_and\_Kohn-Sham\_definition\_C}
\fancyhead[R]{\small\thepage}
\renewcommand{\headrulewidth}{0.4pt}
% 16x16 matrices: tighter columns so that a full matrix fits the text width
\setlength{\arraycolsep}{3.2pt}
% The documents of this project contain long tagged symbols, file paths and
% identifiers that cannot be hyphenated.  A paragraph is therefore allowed to
% have looser lines before TeX gives up and lets a line stick out of the margin.
\AtBeginDocument{\tolerance=2000 \emergencystretch=5em \hbadness=2000}
\ifLuaTeX
  \usepackage{selnolig}  % disable illegal ligatures
\fi
\IfFileExists{bookmark.sty}{\usepackage{bookmark}}{\usepackage{hyperref}}
\IfFileExists{xurl.sty}{\usepackage{xurl}}{} % add URL line breaks if available
\urlstyle{same}
\hypersetup{
  pdftitle={The task, the rules, what was done, what was not, and the decisions that are the author's},
  pdfauthor={Prepared for Patrick Lee Nash, Ph.D.},
  colorlinks=true,
  linkcolor={blue},
  filecolor={Maroon},
  citecolor={blue},
  urlcolor={blue},
  pdfcreator={LaTeX via pandoc}}

\title{The task, the rules, what was done, what was not, and the
decisions that are the author's}
\usepackage{etoolbox}
\makeatletter
\providecommand{\subtitle}[1]{% add subtitle to \maketitle
  \apptocmd{\@title}{\par {\large #1 \par}}{}{}
}
\makeatother
\subtitle{Provenance document \texttt{P0} of PROJECT spacetimes,
Revision\_Deflation\_CJT\_and\_Kohn-Sham\_definition\_C}
\author{Prepared for Patrick Lee Nash, Ph.D.}
\date{November 13, 2026}

\begin{document}
\maketitle

\hypertarget{what-this-document-is}{%
\section{1. What this document is}\label{what-this-document-is}}

This document is the account that the task specification requires of the
work in the folder
\texttt{Revision\_Deflation\_CJT\_and\_Kohn-Sham\_definition\_C}. It
states

\begin{itemize}
\tightlist
\item
  the task, in the words of this provenance, with the author's
  hypotheses in his formulas, the instruction that created this folder,
  and how that instruction was read (Section 2);
\item
  each of his hard rules and how it is met, with the place where that
  can be verified (Section 3);
\item
  what the folder contains (Section 4) and what was done, in order
  (Section 5);
\item
  the results in brief: what the definition by form C changes (Section
  6);
\item
  every \textbf{deviation} from the instructions, with its reason
  (Section 7);
\item
  the \textbf{decisions that are the author's} to take or to confirm
  (Section 8), and the findings that concern his own documents and the
  earlier folders of the project (Section 9);
\item
  what has \textbf{not} been done (Section 10);
\item
  the inputs (Section 11) and the \textbf{log of errors} made during the
  work (Section 12).
\end{itemize}

The four provenance documents of the folder are called \texttt{P0} (this
one), \texttt{P1} (the symbols), \texttt{P2} (the calculations) and
\texttt{P3} (the two forms of the self-interaction), after the beginning
of their file names. Every symbol that occurs here is explained in
Section 13. Nothing in this document has to be looked up elsewhere;
where it names another document of the folder, it does so to say what
that document contains, not to define anything used here.

\hypertarget{the-task}{%
\section{2. The task}\label{the-task}}

\hypertarget{the-hypotheses}{%
\subsection{2.1 The hypotheses}\label{the-hypotheses}}

PROJECT spacetimes studies a sixteen-component Dirac field on two
eight-dimensional spacetimes with the coordinates \(x0,\dots,x7\):
\(x1,x2,x3\) are ordinary space, \(x4\) is the ordinary time, \(x0\) is
a hidden space-like coordinate, and \(x5,x6,x7\) are extra coordinates,
time-like on the \textbf{4+4} spacetime and space-like on the
\textbf{7+1} spacetime.

The author's hypotheses for this sub-project are:

\begin{enumerate}
\def\labelenumi{\arabic{enumi}.}
\tightlist
\item
  \textbf{Deflation.} Each of \(x5\), \(x6\), \(x7\) deflates
  exponentially. Flat spacetime is replaced by a pseudo-Riemannian space
  with a diagonal metric whose components along these coordinates are,
  with three positive rates \(\mathrm{g55},\mathrm{g66},\mathrm{g77}\)
  and positive constants that may differ,
  \[\text{4+4:}\quad g^{(4+4)}_{tt}=-\,\mathrm{const}_t^{\,2}\,e^{-2\,\mathrm{rate}_t\,x4},\qquad
  \text{7+1:}\quad g^{(7+1)}_{tt}=+\,\mathrm{const}_t^{\,2}\,e^{-2\,\mathrm{rate}_t\,x4},\qquad t=5,6,7,\]
  where \(\mathrm{rate}_5=\mathrm{g55}\),
  \(\mathrm{rate}_6=\mathrm{g66}\), \(\mathrm{rate}_7=\mathrm{g77}\).
\item
  \textbf{The Dirac Lagrangians contain a mass term and a
  self-interaction.} The self-interaction is the quartic term of the
  author's document on the 7+1 problem,
  \[\mathcal L_{\rm self}=-\tfrac12\sum_{r=0}^{8}g_r\,I_r,\qquad
  I_r=\sum_{|K|=r}\big(\psi^\dagger\mathrm{Adj}\,\Theta_K\,\psi\big)\big(\psi^\dagger\mathrm{Adj}\,\Theta^K\,\psi\big),\]
  with nine couplings \(g_r\), the 256 products \(\Theta^K\) of Dirac
  matrices and the matrix \(\mathrm{Adj}\) of the adjoint (Section 13).
\item
  \textbf{Two methods.} The effective potential for composite operators
  of Cornwall, Jackiw and Tomboulis (CJT) and the Kohn--Sham ansatz are
  both to be employed: to calculate additional predictions of the two
  models, and as a check on the existing computations. The author's
  solution of the 7+1 problem is to be carried over to the 4+4 problem.
\end{enumerate}

\hypertarget{the-point-at-which-the-71-construction-does-not-carry-over-and-the-authors-decision}{%
\subsection{2.2 The point at which the 7+1 construction does not carry
over, and the author's
decision}\label{the-point-at-which-the-71-construction-does-not-carry-over-and-the-authors-decision}}

On the 7+1 spacetime the adjoint that makes the invariants \(I_r\)
covariant is the adjoint that makes the energy positive:
\(\mathrm{Adj}=\beta=\gamma^4\), \(\bar\psi=\psi^\dagger\beta\). On the
4+4 spacetime the two are different matrices. The covariant adjoint is
\(\tilde\psi=\psi^\dagger C^-\) with
\(C^-=\gamma^0\gamma^1\gamma^2\gamma^3\), and the covariant Lagrangian
\(\tilde\psi(\gamma^\mu\nabla_\mu-m)\psi\) is indefinite: for fields
that do not depend on \(x5,x6,x7\), the only ones with a unitary
evolution (the \emph{quantized sector}), it is

\[\tilde\psi\,(\gamma^A\partial_A-m)\,\psi=-\psi^\dagger\,W_{1644}\,\big(i\partial_4-\mathsf h\big)\,\psi,\qquad
W_{1644}=i\,C^-\gamma^4,\quad\mathsf h=\sum_{j=0}^{3}\alpha^j(-i\partial_j)+\beta\,m,\]

with \(\beta=-i\gamma^4\), \(\alpha^j=-\gamma^4\gamma^j\): the ordinary
Dirac Lagrangian on the half of the field with \(W_{1644}=-1\) and minus
it on the half with \(W_{1644}=+1\). The quantum theory of the free
field is defined by the positive Lagrangian
\(\psi^\dagger(i\partial_4-\mathsf h)\psi\) for all sixteen components.
With a self-interaction the quartic term can then be completed in two
ways:

\begin{itemize}
\tightlist
\item
  \textbf{form C, the covariant form:} the invariants of the covariant
  Lagrangian, \(\mathrm{Adj}=C^-\),
  \(I_r^{\rm C}=\sum_{|K|=r}\mathrm{sgn}_K\,(\tilde\psi\,\Theta^K\psi)^2\);
\item
  \textbf{form Q, the quantized-sector form:} the invariants built with
  the adjoint of the positive Lagrangian, \(\mathrm{Adj}=\beta\), and
  with the matrices \(\gamma_{44q}^A=-i\gamma^A\) (\(A\le4\)),
  \(\gamma_{44q}^A=\gamma^A\) (\(A\ge5\)):
  \(I_r^{\rm Q}=\sum_{|K|=r}\mathrm{sgn}_K^Q\,(\bar\psi\,\Theta_Q^K\psi)^2\).
\end{itemize}

They are related exactly,
\(I_r^{\rm C}=\sum_{|K|=r}\mathrm{sgn}_K^Q\,(\psi^\dagger W_{1644}\,\beta\,\Theta_Q^K\psi)^2\):
form C is form Q with the matrix \(W_{1644}\) inserted in every
bilinear.

The folder \texttt{Revision\_Deflation\_CJT\_and\_Kohn-Sham} of the
repository adopted form Q, reported form C beside it for the couplings,
for uniform matter and for pairs at rest, and put the choice before the
author as his decision (its decision D1). \textbf{On 13 November 2026
the author decided: in a new folder form C is to be the definition; both
options are to be recorded correctly and completely; and the options are
to be explained and discussed in the documents and in the provenance
files.} That instruction is the task of the present folder.

\hypertarget{how-the-instruction-was-read}{%
\subsection{2.3 How the instruction was
read}\label{how-the-instruction-was-read}}

\begin{longtable}[]{@{}
  >{\raggedright\arraybackslash}p{(\columnwidth - 4\tabcolsep) * \real{0.1961}}
  >{\raggedright\arraybackslash}p{(\columnwidth - 4\tabcolsep) * \real{0.5686}}
  >{\raggedright\arraybackslash}p{(\columnwidth - 4\tabcolsep) * \real{0.2353}}@{}}
\toprule\noalign{}
\begin{minipage}[b]{\linewidth}\raggedright
the instruction
\end{minipage} & \begin{minipage}[b]{\linewidth}\raggedright
what was done
\end{minipage} & \begin{minipage}[b]{\linewidth}\raggedright
where
\end{minipage} \\
\midrule\noalign{}
\endhead
\bottomrule\noalign{}
\endlastfoot
a new folder & The folder
\texttt{Revision\_Deflation\_CJT\_and\_Kohn-Sham\_definition\_C}, with
its own program (\texttt{deflation\_cjt\_ks\_c}), data, figures and
documents. The folder \texttt{Revision\_Deflation\_CJT\_and\_Kohn-Sham}
was not changed in any file; nor were the folder
\texttt{Revision\_Deflation} and the flat-spacetime part of the project.
& the repository \\
form C is the definition & In every document of the folder the 4+4 model
is defined by
\(\mathcal L_{\rm kin}+\mathcal L_{\rm mass}+\mathcal L_{\rm self}^{\rm C}\):
the positive kinetic and mass terms of the quantized sector with the
quartic term of the covariant Lagrangian. Every result of the 4+4 model
that depends on the self-interaction was computed again with form C: the
mean-field couplings, uniform matter (now of two halves of the field),
the effective potential and the vacuum, the pair couplings and the gap,
the symmetries, the matter as a source of gravity, the thresholds with
the Vorton-Scalar. & manuscripts; ``dirac matrices deflating'', Parts 5
and 6; \texttt{P3} \\
``C'' with which kinetic term & The covariant kinetic term is indefinite
and defines no quantum theory with states of positive norm (Section
2.2). Form C therefore means the covariant quartic term together with
the positive kinetic term. The third possibility, the covariant kinetic
term with the covariant quartic term, is recorded as a classical field
equation,
\(i\partial_4\psi=\mathsf h\psi+W_{1644}\,\partial\mathcal L_{\rm self}^{\rm C}/\partial\psi^\dagger\);
nothing is computed with it (Section 7, V8). & the same \\
both options recorded correctly and completely & Every calculation that
involves the self-interaction on 4+4 is carried out for form C and for
form Q with the same program, the same couplings of the Lagrangian and
the same routines, and every table of such results prints both. The
exact relation between the forms is verified for all 256 matrices in
exact arithmetic. An independent script in another language repeats the
comparison of the two forms, and deliberate errors introduced into it
are detected. & \texttt{P3}; \texttt{P2}, Sections 5.1, 5.3, 5.4, 5.13,
5.14 \\
explain and discuss the options in the documents and in the provenance
files & A provenance document of its own, \texttt{P3}, written for a
reader who knows neither form; a section ``The two forms of the
self-interaction: discussion'' in the 4+4 manuscript; Section 6.7 of
``dirac matrices deflating''; Sections 2, 6, 8 and 9 of this document;
Section 3.2 of \texttt{P1} for the symbols; Section 2 of \texttt{P2} for
the defining equations. Each discussion states what each form is, what
speaks for it and against it, what does not distinguish the forms, and
what could decide. & the documents named \\
\end{longtable}

\hypertarget{the-criticisms-that-the-earlier-task-states-and-what-was-done-about-each}{%
\subsection{2.4 The criticisms that the earlier task states, and what
was done about
each}\label{the-criticisms-that-the-earlier-task-states-and-what-was-done-about-each}}

The task specification of the sub-project (of which this folder is the
third revision) states three criticisms of the first revision. They hold
for this folder as they did for its predecessor.

\begin{longtable}[]{@{}
  >{\raggedright\arraybackslash}p{(\columnwidth - 2\tabcolsep) * \real{0.5000}}
  >{\raggedright\arraybackslash}p{(\columnwidth - 2\tabcolsep) * \real{0.5000}}@{}}
\toprule\noalign{}
\begin{minipage}[b]{\linewidth}\raggedright
criticism
\end{minipage} & \begin{minipage}[b]{\linewidth}\raggedright
what was done
\end{minipage} \\
\midrule\noalign{}
\endhead
\bottomrule\noalign{}
\endlastfoot
The manuscripts are not useful for students: they send the reader to
other manuscripts, documents, provenances or symbol files. & Each
manuscript contains every symbol, definition and equation that it uses,
with appendices for the symbols, the printed Dirac matrices, the
curvature, and the programs and checks. A script searches every document
for wordings that send the reader elsewhere
(\texttt{check\_self\_contained.py}); none is found. \\
The manuscripts do not define the Dirac Lagrangians and Dirac equations
correctly, because the self-interaction is missing. & The Lagrangians
are \(\mathcal L_{\rm kin}+\mathcal L_{\rm mass}+\mathcal L_{\rm self}\)
on both spacetimes, on 4+4 in both forms, and the field equations carry
the cubic term; manuscripts, Secs. III and IV, and the document ``dirac
matrices deflating'', Parts 3 to 6. \\
The same symbol is used for different objects. & A MASTER registry of
246 symbols with one meaning each, the list of the collisions found and
their resolution, the symbols of the two forms, and eight machine
checks: document \texttt{P1}. All documents of the folder use it. \\
\end{longtable}

\hypertarget{the-hard-rules-and-how-each-is-met}{%
\section{3. The hard rules, and how each is
met}\label{the-hard-rules-and-how-each-is-met}}

\begin{longtable}[]{@{}
  >{\raggedright\arraybackslash}p{(\columnwidth - 4\tabcolsep) * \real{0.3333}}
  >{\raggedright\arraybackslash}p{(\columnwidth - 4\tabcolsep) * \real{0.3333}}
  >{\raggedright\arraybackslash}p{(\columnwidth - 4\tabcolsep) * \real{0.3333}}@{}}
\toprule\noalign{}
\begin{minipage}[b]{\linewidth}\raggedright
rule
\end{minipage} & \begin{minipage}[b]{\linewidth}\raggedright
how it is met
\end{minipage} & \begin{minipage}[b]{\linewidth}\raggedright
where to verify
\end{minipage} \\
\midrule\noalign{}
\endhead
\bottomrule\noalign{}
\endlastfoot
0. The audience is a group of clever physics Ph.D.~students who do not
know the solution; they must be able to learn from it and to teach it. &
Every term and symbol is defined in words before it is used; derivations
are given in steps that can be repeated by hand, and the text says how;
each result names the check that verifies it. & any section of a
manuscript, of ``dirac matrices deflating'' or of \texttt{P3} \\
1. All documentation, symbol and provenance files exist as .md, .tex and
.pdf, complete, precise and correct. & Each of the eight documents of
Section 4 exists in the three forms; the .tex and .pdf are built from
the .md (for the manuscripts: the .md from the .tex) by scripts, and the
build fails on an undefined reference or on a box wider than its space.
& the files; \texttt{code/scripts/build\_docs.sh},
\texttt{build\_manuscripts.py} \\
2. Push to the repository as soon as an object is updated. & Every
commit is carried at once to the author's folder and fast-forwarded
there; the push itself is done from the author's computer (Section 7,
V3). At the beginning of the work on this folder the rule was broken
(Section 7, V9). & \texttt{git\ log} \\
3., 4. Working directory and repository. & The folder
\texttt{Revision\_Deflation\_CJT\_and\_Kohn-Sham\_definition\_C} of the
repository \texttt{https://github.com/once-ere/spacetimes}, in the
working directory named by the rule. & \\
5. No symbol with two meanings; a MASTER provenance of the symbols,
promulgated and used. & Document \texttt{P1} and the table
\texttt{provenance/symbols\_master.csv}; the checks MS1 to MS8 (8 of 8
pass) verify the registry, the documents and the columns of all data
files against it. The first version of \texttt{P3} broke the rule
(Section 12, N-E9). & \texttt{P1}, Section 5;
\texttt{data/symbol\_checks.txt} \\
6. Source code for every equation, result, figure and prediction; a
figure for each prediction; commented code, instructions and provenance
for every computed value; Rust preferred; the code executed. & One Rust
program (36 source files, 17774 lines, 77 unit tests, 19 commands)
computes everything; 38 figures are drawn from its data files; document
\texttt{P2} gives, for every table and figure, the command, the source
files and the data files, and prints the complete output of every check
(310 checks of the program, all passing). Every command was executed in
this folder; the reports in \texttt{data} are the outputs, and a check
verifies that every data file was written here. Python is used for three
purposes only (Section 7, V2). & \texttt{P2} \\
7. Never send the reader to another manuscript, document, provenance or
symbol file; every equation, symbol and definition in the file itself. &
Each document defines what it uses; \texttt{P3}, for example, contains
the complete registry of symbols and the equations of gravity that its
section on gravity needs. Program output printed in the documents uses
the expression ``the main project'' for the flat-spacetime part of
PROJECT spacetimes, whose program is used as a library and whose data
files several checks compare against; the documents that print such
output say so. & \texttt{data/self\_contained\_checks.txt} \\
8. Follow the instructions exactly, or explain in a provenance why not,
and pause. No shortcuts. & Sections 7 to 10 and 12 of this document.
& \\
\end{longtable}

\hypertarget{what-the-folder-contains}{%
\section{4. What the folder contains}\label{what-the-folder-contains}}

\begin{longtable}[]{@{}
  >{\raggedright\arraybackslash}p{(\columnwidth - 2\tabcolsep) * \real{0.5000}}
  >{\raggedright\arraybackslash}p{(\columnwidth - 2\tabcolsep) * \real{0.5000}}@{}}
\toprule\noalign{}
\begin{minipage}[b]{\linewidth}\raggedright
document (each as .md, .tex, .pdf)
\end{minipage} & \begin{minipage}[b]{\linewidth}\raggedright
content
\end{minipage} \\
\midrule\noalign{}
\endhead
\bottomrule\noalign{}
\endlastfoot
\texttt{manuscript/16\ Component\ Dirac\ Fermion\ Condensate\ Singlet\ Superfluid\ on\ 4+4\ Deflating\ Spacetime}
& the manuscript for the 4+4 spacetime, in the style of Physical Review
D; form C is the definition, every result of form Q is given beside it,
and Sec. XVIII D discusses the two \\
\texttt{manuscript/16\ Component\ Dirac\ Fermion\ Condensate\ Singlet\ Superfluid\ on\ 7+1\ Deflating\ Spacetime}
& the manuscript for the 7+1 spacetime, on which the self-interaction
has one form \\
\texttt{dirac\ matrices\ deflating} & geometry, the sixteen Dirac
matrices printed in full, spin connection, the Lagrangians with mass
term and self-interaction, the field equations, the two forms of the 4+4
model and their algebra, their consequences in numbers and their
discussion, solutions mode by mode; complete reports of the checks \\
\texttt{provenance/P0\_task\_rules\_deviations\_decisions} & this
document \\
\texttt{provenance/P1\_master\_symbols\_spacetimes} & the MASTER
registry of symbols, the collisions and their resolution, the symbols of
the two forms, the dictionary to the author's 7+1 document, the machine
checks \\
\texttt{provenance/P2\_calculations\_code\_and\_checks} & every
calculation: defining equations, commands, instructions, data files,
figures, complete output of all checks, source files \\
\texttt{provenance/P3\_two\_forms\_of\_the\_self\_interaction} & the two
forms from the beginning: the indefinite covariant Lagrangian, both
definitions with their field equations, the exact relation, uniform
matter, vacuum, pairs, symmetries, gravity, what has been established
and what has not, the discussion \\
\texttt{README} & the list of contents and the instructions for building
everything \\
\end{longtable}

\begin{longtable}[]{@{}
  >{\raggedright\arraybackslash}p{(\columnwidth - 2\tabcolsep) * \real{0.5000}}
  >{\raggedright\arraybackslash}p{(\columnwidth - 2\tabcolsep) * \real{0.5000}}@{}}
\toprule\noalign{}
\begin{minipage}[b]{\linewidth}\raggedright
other content
\end{minipage} & \begin{minipage}[b]{\linewidth}\raggedright
\end{minipage} \\
\midrule\noalign{}
\endhead
\bottomrule\noalign{}
\endlastfoot
\texttt{code/deflation-cjt-ks-c-rs} & the Rust program
\texttt{deflation\_cjt\_ks\_c} \\
\texttt{code/scripts} & independent verification in Python (two scripts
and a test of them with deliberate errors); figures; the scripts that
build and check the documents \\
\texttt{data} & 126 data files written by the program, and the reports
of all checks \\
\texttt{figures} & 38 figures, each as .pdf and .png \\
\texttt{docs\_src} & the sources of all documents \\
\end{longtable}

The two manuscripts are built from one source
(\texttt{docs\_src/manuscript}): the text that the two share is written
once, and passages that differ are marked for one spacetime or the
other, so that the manuscripts cannot drift apart.

\hypertarget{what-was-done-in-order}{%
\section{5. What was done, in order}\label{what-was-done-in-order}}

The folder began as a copy of the tracked files of the folder
\texttt{Revision\_Deflation\_CJT\_and\_Kohn-Sham}. Nothing of that copy
was committed before it had been produced again in this folder: every
command of the program was executed here, every figure drawn here, every
document rebuilt here.

\begin{longtable}[]{@{}
  >{\raggedright\arraybackslash}p{(\columnwidth - 2\tabcolsep) * \real{0.5000}}
  >{\raggedright\arraybackslash}p{(\columnwidth - 2\tabcolsep) * \real{0.5000}}@{}}
\toprule\noalign{}
\begin{minipage}[b]{\linewidth}\raggedright
commit
\end{minipage} & \begin{minipage}[b]{\linewidth}\raggedright
content
\end{minipage} \\
\midrule\noalign{}
\endhead
\bottomrule\noalign{}
\endlastfoot
\texttt{d392864} & the program with form C as the definition: mean-field
couplings of four densities for both forms, matter of the two halves,
the CJT potential of two halves, the vacuum, both forms in the pair
channels, matter of both forms as a source of gravity \\
\texttt{6d654e8} & correction of the search for the state of lowest
energy per particle (Section 12, N-E2) \\
\texttt{782315c}, \texttt{85ca51a} & data of the fourteen commands that
do not involve the self-interaction, executed in this folder; identical
in content to those of the folder with form Q \\
\texttt{8277767} & data of the commands \texttt{selfcoupling},
\texttt{kohnsham}, \texttt{pairing}, \texttt{gravity} \\
\texttt{69d03e8} & the ultraviolet criterion of the gap equation, and
the upper limit of its momentum integral raised (Section 9, F6 and
F7) \\
\texttt{21f288b} & data of the command \texttt{forms}: the two forms
side by side (corrections N-E4 to N-E6) \\
\texttt{9f8a206} & \texttt{P3}, first version; the independent
verification of the two forms; all figures \\
\texttt{2ef3704} & the symbols of the two forms entered in the MASTER
registry; \texttt{P3} rewritten with them; the names of the program and
of the data columns changed accordingly (N-E9, N-E10) \\
\texttt{0cd79e2} & ``dirac matrices deflating'' with form C as the
definition \\
\texttt{2b81eb2} & the test of the independent scripts with deliberate
errors \\
\texttt{25d5665}, \texttt{a824eab} & the two manuscripts; corrections
after reading the typeset pages (N-E15, N-E16) \\
the commits that follow & \texttt{P1}, \texttt{P2}, \texttt{README},
this document; the check that every data file was written in this
folder \\
\end{longtable}

\hypertarget{the-results-in-brief-what-the-definition-by-form-c-changes}{%
\section{6. The results in brief: what the definition by form C
changes}\label{the-results-in-brief-what-the-definition-by-form-c-changes}}

Each statement is verified by the checks named in the manuscripts, in
\texttt{P3} and in \texttt{P2}. Numbers are in units of the fermion mass
\(m\) unless said otherwise. The sets of couplings are \textbf{set A}, a
scalar attraction alone, \(g_0=-g\), and \textbf{set B}, \(g_0=-g\) with
a vector repulsion one tenth as large, \(g_1=+0.1\,g\), with a strength
\(g>0\) that grows as \(e^{\tau_d}\) along the deflation clock
\(\tau_d\).

\textbf{What is the same in both forms.} Everything that does not
involve the contact self-interaction: the geometry, the Dirac operator,
the solutions mode by mode, the exchange of the Vorton-Scalar (whose
Yukawa coupling contains \(\bar\psi\) in both forms), the cosmological
part. The reduction of the nine couplings to four effective couplings
\(g_0^{\rm eff}\), \(g_1^{\rm eff}\), \(g_7^{\rm eff}\),
\(g_8^{\rm eff}\), with the same Fierz matrix. The interaction between
two particles of the same half of the field. The equivalence of the
Kohn--Sham method and the CJT effective potential, which agree for the
two halves in both forms to \(4.9\times10^{-11}\) on 80 states. On the
7+1 spacetime there is one form.

\textbf{The mean field.} The energy density of the self-interaction in
uniform matter is, in both forms, a sum of four squares, of the total
number density and scalar density and of the differences of these
densities between the two halves; the four coefficients are the same
four combinations of the effective couplings, and the two forms exchange
the coefficients of the totals with those of the differences. For set A
the attraction \(\tfrac{15}{16}g_0\) acts on the total scalar density in
form Q and on the difference of the scalar densities in form C. Between
two particles of the same half the coupling is \(\tfrac78g_0\) in both;
between particles of different halves it is \(g_0\) in form Q and
\(-g_0\) in form C.

\textbf{Matter.} In form C matter polarizes: the state of lowest energy
is polarized in 48 of the 48 cases examined in form C and in 0 of them
in form Q. Self-bound matter of form C has one half of the field filled,
that of form Q both halves equally. The strength \(g\,m^{D-1}\) at which
a gas becomes self-bound matter is, in \(D=4\) space dimensions with
momentum, 224.74 (form C) against 82.11 (form Q) for set A and 332.29
against 121.09 for set B; in \(D=3\), 24.572 against 9.809 and 31.167
against 12.382: form C needs 2.51 to 2.74 times the strength and binds
later on the clock by 0.92 to 1.01. A gas of fixed density polarizes
before it binds.

\textbf{Vacuum.} With the filled states of negative energy below a
cutoff \(\Lambda\) a field without mass term acquires a mass above the
same critical strength in both forms (31.583 in units of
\(\Lambda^{-3}\) for \(D=4\), set A), and of the same magnitude; in form
Q the two halves get the same mass, in form C masses of opposite sign,
which is the mass term \(\tilde\psi\psi\) of the covariant Lagrangian.
With a bare mass form Q passes smoothly through the critical strength,
and form C has a sharp transition at which the symmetry between the
halves breaks.

\textbf{Pairs.} A fermion pair along one of the space-like axes
\(x0,\dots,x3\) consists of one particle of each half, a pair along one
of the additional times \(x5,x6,x7\) of two particles of the same half.
In form C the contact couplings therefore enter the gap equation of the
pair along the hidden axis with the sign opposite to that of form Q
(amplitudes compared at random momenta, \(2.2\times10^{-15}\)): a scalar
contact attraction lowers its gap, a vector repulsion raises it, and the
four pairs along \(x0,x5,x6,x7\), which are degenerate in form Q, are
not. At the clock reading 2.04 of the example (\(m_s=m\), Fermi momentum
\(0.5\,m\)) the exchange of the Vorton-Scalar alone gives the gap 0.0247
of the Fermi kinetic energy; the scalar contact attraction of the
example multiplies it by 12.1 in form Q and removes it in the pair along
the hidden axis of form C.

\textbf{Symmetries.} The self-interaction is chirally invariant if
\(g_1^{\rm eff}-g_7^{\rm eff}=0\) in form C and if
\(g_0^{\rm eff}+g_8^{\rm eff}=0\) in form Q. Every self-interaction of
form C that does not vanish breaks the internal symmetry
\(SO(4)\times U(1)_E\) of fields in ordinary space, on which the
hypothesis about the Higgs field rests, down to the rotations of
\(x5,x6,x7\); form Q keeps \(SO(4)\) for all couplings.

\textbf{Gravity.} With self-bound matter of set B at three times its
binding threshold as the source of the tuned Einstein--Lovelock
equations (the gravity in which the undisturbed deflation is a solution
without matter), and with the smaller of the two amounts of matter
examined, the rate of deflation at the end of the evolution is 34.1
times the initial rate for the matter of form C and 15.6 times for that
of form Q; the conclusions (matter accelerates the deflation; the
undisturbed deflation is unstable) are the same.

\textbf{A result that concerns both forms and both spacetimes: the
ultraviolet criterion.} The gap equation has a solution that does not
depend on the upper limit of its momentum integral only if the
ultraviolet strength of its kernel, \(U_{\rm uv}\), is below \(4\pi^2\)
(Section 9, F6 and F7).

\hypertarget{deviations-from-the-instructions-and-the-reasons}{%
\section{7. Deviations from the instructions, and the
reasons}\label{deviations-from-the-instructions-and-the-reasons}}

\textbf{V1. The five supplements of the first revision are not repeated
as separate documents.} Hard Rule 7 forbids a document to send its
reader to another, and a supplement is a document to which a manuscript
sends its reader. The content of the supplements is therefore inside
each manuscript:

\begin{longtable}[]{@{}
  >{\raggedright\arraybackslash}p{(\columnwidth - 2\tabcolsep) * \real{0.5000}}
  >{\raggedright\arraybackslash}p{(\columnwidth - 2\tabcolsep) * \real{0.5000}}@{}}
\toprule\noalign{}
\begin{minipage}[b]{\linewidth}\raggedright
supplement of the first revision
\end{minipage} & \begin{minipage}[b]{\linewidth}\raggedright
where its content is in each manuscript of this revision
\end{minipage} \\
\midrule\noalign{}
\endhead
\bottomrule\noalign{}
\endlastfoot
Vorton-Scalar and superfluid & Secs. VI, VIII, X, XI, XIV \\
effective action along the clock & Sec. XII \\
cosmology & Sec. XVI A to D, F, I to K \\
microwave background & Sec. XVI E, G, H \\
Higgs hypothesis & Sec. XVII \\
\end{longtable}

The derivations are in the sections named. Of the tables of the
supplements, the sections print those that carry the argument and quote
numbers from the others; the others are printed in full in the appendix
``Complete tables of the calculations'' of each manuscript, each with a
caption that defines its columns.

\textbf{V2. Python is used beside Rust}, for three purposes: (i) the
independent verification, which must not share code with the program
that it verifies and therefore uses another language and other libraries
(computer algebra, a general-purpose integrator, other special
functions); (ii) the figures; (iii) the scripts that assemble, typeset
and check the documents. No result of a manuscript comes from Python
alone. Mathematica was not available.

\textbf{V3. The push to GitHub is done from the author's computer.} The
work is carried out in a workspace without the author's credentials: a
push from the workspace is refused by the server, and none is possible
from the part of his computer that the work can reach. Each commit is
transferred at once to his folder as a git bundle and fast-forwarded
there; the folder is then ahead of the remote until he runs
\texttt{git\ push\ origin\ main} in it.

\textbf{V4. The test with deliberate errors does not cover the whole of
the first independent script.} It introduces 45 errors (29 in the script
for the self-interaction, the Kohn-Sham and CJT calculations, pairing
and gravity, 16 in the cosmological part of the other script); all are
detected. The part of the first script that verifies the geometry, the
Dirac operators and the modes has no such test.

\textbf{V5. The titles of the manuscripts} are those of the first
revision; the task names no titles.

\textbf{V6. An external Einstein--Boltzmann code was not used as a
cross-check.} The spectra of the microwave background are computed with
the program of this project, whose version for static extra coordinates
had been validated in the flat-spacetime part of the project; no
external code was built or run in this sub-project.

\textbf{V7. The 7+1 manuscript is built in this folder too}, although
the instruction concerns the 4+4 model. One source builds both
manuscripts, and the finding on the upper limit of the gap equation
(Section 9, F6) concerns both: the 7+1 manuscript of this folder
contains the ultraviolet criterion and gaps computed with the higher
limit, which differ in the last digits from those of the folder with
form Q. In everything else it has the content of its predecessor.

\textbf{V8. Nothing is computed with the third option}, the covariant
quartic term with the covariant (indefinite) kinetic term. Its field
equation is derived and recorded in every document that defines the
forms. It conserves \(\psi^\dagger W_{1644}\psi\) and not
\(\psi^\dagger\psi\), and it is not the equation of a Hamiltonian with
states of positive norm when both halves of the field are occupied; no
quantum theory of particles, and therefore no matter, vacuum or pairing,
can be computed from it with the methods of this project.

\textbf{V9. Hard Rule 2 was broken at the beginning of the work on this
folder.} From the instruction (14:06, Pacific time) to the first commit
of the folder (15:34) nothing was committed or carried to the author's
folder, while the program was being extended and the long commands were
running. The author pointed it out. Since then every step ends with a
commit that is transferred at once (Section 12, N-E1).

\textbf{V10. The top-level \texttt{README.tex} and \texttt{README.pdf}
of the repository were not rebuilt} after the author changed a word in
the top-level \texttt{README.md} (his commit \texttt{43e4525},
``Licence'' to ``License''): the flat-spacetime part of the project, to
which these files belong, is paused by his instruction. His change was
merged; the three files differ in that one word until that part is taken
up again.

\hypertarget{decisions-that-are-the-authors}{%
\section{8. Decisions that are the
author's}\label{decisions-that-are-the-authors}}

\textbf{D1. The form of the self-interaction on the 4+4 spacetime:
decided.} Form C, by the author's decision of 13 November 2026 (Section
2.2). Two matters connected with it remain his:

\begin{itemize}
\tightlist
\item
  \emph{Which document stands for the 4+4 model.} The repository now
  holds two 4+4 manuscripts with the same title, one in the folder with
  form Q as the definition and one in this folder. If this folder
  supersedes the other, the other should say so; it was not changed.
\item
  \emph{The sets of couplings.} The numbers of this folder are given for
  the sets A and B. Set B has the ratio of the vector to the scalar
  coupling of the author's reference example. The formulas hold for all
  couplings; other sets are a matter of running the command
  \texttt{forms} with them.
\end{itemize}

\textbf{D2. The bare mass.} With the Vorton-Scalar, a bare mass \(m_D\)
breaks the chiral symmetry explicitly and gives the phase mode a mass.
Results that rest on a massless phase mode (the stepped fluid of the
cosmological part) presuppose \(m_D=0\) or a mass of the phase mode
below the temperatures concerned. Which is meant is the author's
decision.

\textbf{D3. The couplings of the reference example.} See Section 9, F2:
are \(-10^5\) and \(+10^4\) the couplings \(g_0\), \(g_1\) of the
Lagrangian or the couplings \(G_S\), \(G_V\) of the mean-field energy?

\textbf{D4. The symbols of the metric.} The metric of this sub-project
is written \(g^{(4+4)}\), \(g^{(7+1)}\) in the author's sign convention
(space-like components positive), and the rates keep his names
\(\mathrm{g55},\mathrm{g66},\mathrm{g77}\), upright. The names of the
first revision for a metric of the other sign convention are retired.

\textbf{D5. Open from the first revision:} the assumption about the
perturbations of ordinary space and its closures; the nature of whatever
sustains the deflation; the choice between Einstein gravity with such an
agent and Einstein--Gauss--Bonnet gravity.

\textbf{D6. The earlier folders and the upper limit of the gap
equation.} The folders \texttt{Revision\_Deflation\_CJT\_and\_Kohn-Sham}
and \texttt{Revision\_Deflation} and the flat-spacetime part of the
project solve the gap equation with an upper limit of the momentum
integral that is too low (Section 9, F6). They were not changed: the
first by the instruction that created this folder, the others because
they are finished or paused. Whether they are to be computed again with
the higher limit is the author's decision.

\textbf{D7. The hypothesis about the Higgs field under form C.} The
pattern on which the hypothesis rests, a doublet under a hidden
\(SO(4)\) that the condensate along the hidden axis breaks to \(SO(3)\),
is present in the model as now defined only if there is no contact
self-interaction at all (Section 9, F5). Whether the hypothesis is to be
kept for the model without contact self-interaction, tied to form Q, or
given up, is the author's decision.

\hypertarget{findings-that-concern-the-authors-own-documents-and-the-earlier-folders}{%
\section{9. Findings that concern the author's own documents and the
earlier
folders}\label{findings-that-concern-the-authors-own-documents-and-the-earlier-folders}}

\textbf{F1. Only four of the nine couplings are independent.} The
identity \((1+\mathsf F)I=0\) holds for anticommuting fields with the
Fierz matrix of his document, in both forms. A tensor coupling
\(g_2,\dots,g_6\) is not a new parameter.

\textbf{F2. The couplings of the Kohn--Sham example are not the
Hartree--Fock couplings of \(g_0=-10^5\), \(g_1=+10^4\).} His document
states the Hartree--Fock couplings of these values correctly as
\(-98750\) and \(+20000\) (check \texttt{self-X1}), and its Kohn--Sham
calculation uses \(G_S=-10^5\), \(G_V=+10^4\). With the first pair,
matter on flat 7+1 spacetime is not self-bound: its energy per particle
has only a local minimum above the mass. The numbers of his document are
reproduced with the second pair. If the mean-field couplings are meant
to be the Hartree--Fock couplings, the Lagrangian couplings that give
his example are \(g_0=-105341.6\) and \(g_1=+2484.5\) in the same units.

\textbf{F3. The coupling of the example exceeds the critical coupling of
the vacuum} when the filled states of negative energy are included with
a cutoff above 0.6665: the Kohn--Sham calculation without those states
presupposes a cutoff below that value. The same holds for self-bound
matter of the 4+4 model by contact couplings alone, in both forms and
more so in form C, which needs the stronger couplings: at each binding
threshold the largest cutoff for which the vacuum is stable is below the
mass of the particles.

\textbf{F4. The deflating spacetime is an unstable solution} of the
gravity in which it is a vacuum solution, and in Einstein gravity it
needs a source without energy density and with a tension; Kohn--Sham
matter cannot be that source, in either form.

\textbf{F5. Form C breaks the internal symmetry on which the hypothesis
about the Higgs field rests.} In form Q, and on 7+1, a vector or
axial-vector coupling breaks \(U(1)_E\), the second phase symmetry of
the fields in ordinary space, and every coupling keeps the rotations
\(SO(4)\) of the four extra axes. In form C every self-interaction that
does not vanish breaks the three rotations that mix the hidden axis with
\(x5,x6,x7\) and breaks \(U(1)_E\) (check \texttt{44c-S12}, exact): the
transformations generated by \(\gamma^0\gamma^t\) are boosts of the 4+4
frame, and the unitary transformations generated by the same matrices
are not elements of the frame group, so that a covariant quartic term
has no reason to be invariant under them. With an attractive scalar
contact coupling the pairs along the additional times, not the pair
along the hidden axis, have the larger gap.

\textbf{F6. The gap equation of the earlier calculations was solved with
an upper limit of the momentum integral that is too low.} The folders
named in D6 integrate to 60 times the larger of \(m\) and \(k_F\). The
kernel falls only as the inverse square of the momentum, so that the gap
function has a tail that falls as a power. Measured in this folder with
the limits 60, \(2\times10^3\), \(2\times10^4\) and \(2\times10^5\) on
192 kernels (check \texttt{pc-P7}): with the limit 60 the gaps of the
exchange of the Vorton-Scalar alone are too small by up to 1.08 per cent
at the strongest coupling of the histories, and gaps with attractive
contact couplings by more, the more the nearer the kernel is to the
criterion of F7 (the complete table by classes of the ultraviolet
strength is in the manuscripts and in \texttt{P3}). Quantities derived
from the gap function change with it: in the comparison of the
Ginzburg--Landau coefficients with their weak-coupling formulas the
largest deviation at the third moment is 9.8 per cent with the higher
limit, where the limit 60 had given 7.9 per cent (check
\texttt{def-P3}). This folder uses \(2\times10^4\) throughout, at which
the gaps below the criterion have converged. The conclusions of the
earlier folders are not affected; their last digits are.

\textbf{F7. The gap equation of the model has no solution beyond a
definite coupling.} A solution that does not depend on the upper limit
exists only if the ultraviolet strength of the kernel,

\[U_{\rm uv}=2\,u_S\,g_{VS}^2-s_{\rm ct}\,m_{\rm ct}^2\big[(g_0^{\rm eff}-g_8^{\rm eff})\,u_S+(g_1^{\rm eff}+g_7^{\rm eff})\,u_V\big],
\qquad u_S=\frac{\pi^2}{8}-\frac12,\quad u_V=\frac{3\pi^2}{4}-1,\]

is below \(4\pi^2\); above it the gap grows in proportion to the upper
limit (the pair collapses to short distances, as two relativistic
particles do in a Coulomb potential that is too strong). The exchange of
the two modes of the Vorton-Scalar alone reaches the criterion at
\(g_{VS}^2=26.90\), the clock reading 3.29, beyond the last reading 2.77
of the histories of the project; attractive contact couplings bring it
nearer. No entry of the gap tables of the earlier folders lies above the
criterion; the entries nearest to it, those with the scalar contact
attraction at the last clock reading, are the ones that the limit 60
changes most (F6). In this folder the contact couplings enter the pair
along the hidden axis of form C with the opposite sign, a vector
repulsion then acts as an attraction, and with the vector coupling of
the examples the criterion is reached at the clock reading 2.46; no gap
is given above the criterion. Along the deflation clock every coupling
grows without bound, so that the criterion is reached at a definite
moment in every history: beyond it the model needs physics at short
distances that it does not contain.

\hypertarget{what-has-not-been-done}{%
\section{10. What has not been done}\label{what-has-not-been-done}}

\begin{enumerate}
\def\labelenumi{\arabic{enumi}.}
\tightlist
\item
  \emph{Which form is the right one.} No calculation of this folder
  decides between form C and form Q. The documents set out what speaks
  for each and what could decide: a principle (a covariant theory on the
  full 4+4 spacetime, or a theory of eight equivalent internal states),
  a consistent quantization of the indefinite covariant Lagrangian, or
  an observable consequence.
\item
  \emph{States with coherence between the two halves of the field.} The
  density matrices of this folder are diagonal in \(W_{1644}\). Because
  the interaction connects the halves, states with a coherence between
  them are possible in principle, in both forms; whether one has a lower
  energy has not been examined.
\item
  \emph{Other couplings.} The numbers are for the sets A and B; the
  exact formulas hold for all couplings.
\item
  \emph{The continuation of form Q with a mass that has changed its
  sign.} Where the gap equations of form Q have no solution with
  non-negative masses (108 of 3936 states examined, all with unequal
  halves), the continuation was not followed.
\item
  \emph{The potential of the order parameter with the contact couplings
  of form C}, and \emph{pairing inside polarized matter}: in matter of
  form C with one filled half a pair along the hidden axis, which needs
  one particle of each half, cannot form at all; the gaps are computed
  for particles of the mass \(m\) without regard to the polarization.
\item
  \emph{Matter beyond the mean field; pairing inside dense self-bound
  matter; a self-consistent solution with the created particles acting
  on the metric.}
\item
  \emph{A test with deliberate errors for the geometric part of the
  first independent script} (Section 7, V4); \emph{separate supplements}
  (V1).
\item
  \emph{The earlier folders were not corrected} for the upper limit of
  the gap equation (Section 8, D6), and \emph{the paused stages of the
  flat-spacetime part of the project} were not touched (Section 7, V10).
\end{enumerate}

\hypertarget{inputs}{%
\section{11. Inputs}\label{inputs}}

\begin{longtable}[]{@{}
  >{\raggedright\arraybackslash}p{(\columnwidth - 2\tabcolsep) * \real{0.5000}}
  >{\raggedright\arraybackslash}p{(\columnwidth - 2\tabcolsep) * \real{0.5000}}@{}}
\toprule\noalign{}
\begin{minipage}[b]{\linewidth}\raggedright
input
\end{minipage} & \begin{minipage}[b]{\linewidth}\raggedright
what was taken from it
\end{minipage} \\
\midrule\noalign{}
\endhead
\bottomrule\noalign{}
\endlastfoot
the author's task specification for the sub-project (private; not in the
repository) & the hypotheses and the rules of Sections 2.1 and 3 \\
the author's instruction of 13 November 2026 & the task of this folder
(Section 2.2) \\
the folder \texttt{Revision\_Deflation\_CJT\_and\_Kohn-Sham} of the
repository & the program, the scripts and the sources of the documents
from which this folder started; it was not changed \\
his document \texttt{KOHN\_SHAM\_GRAVITY} (folder \texttt{7+1} on his
computer; its .tex source was read in full) & the definition of
\(\mathcal L_{\rm self}\), of \(\Theta^K\), \(I_r\), \(g_r\); the Fierz
matrix; the Kohn--Sham equations and the reference numbers; the
Einstein--Lovelock equations for two scale factors and the reference
evolution \\
his document \texttt{DIRAC16COMPLEX} (same folder) & the critical
coupling of the vacuum in seven space dimensions \\
the notebooks and papers of the project, as recorded in the
flat-spacetime part & the two sets of real \(16\times16\) Dirac matrices
(files \texttt{data/gamma1644.txt}, \texttt{data/gamma1671.txt} at the
top of the repository); the conventions of the fields \\
the folder \texttt{Revision\_Deflation} (the first revision) & the
calculations that are repeated here; it was not changed \\
\end{longtable}

The author's files were read, not changed, and none of them is copied
into the repository.

\hypertarget{log-of-errors}{%
\section{12. Log of errors}\label{log-of-errors}}

Every error made during the work is recorded here, including errors that
were corrected before anything was committed. Section 12.1 lists the
errors made in the work on this folder. Section 12.2 repeats the log of
the folder \texttt{Revision\_Deflation\_CJT\_and\_Kohn-Sham}, from which
this folder started and whose program and documents it inherited with
their corrections.

\hypertarget{errors-made-in-the-work-on-this-folder}{%
\subsection{12.1 Errors made in the work on this
folder}\label{errors-made-in-the-work-on-this-folder}}

\begin{longtable}[]{@{}
  >{\raggedright\arraybackslash}p{(\columnwidth - 6\tabcolsep) * \real{0.2500}}
  >{\raggedright\arraybackslash}p{(\columnwidth - 6\tabcolsep) * \real{0.2500}}
  >{\raggedright\arraybackslash}p{(\columnwidth - 6\tabcolsep) * \real{0.2500}}
  >{\raggedright\arraybackslash}p{(\columnwidth - 6\tabcolsep) * \real{0.2500}}@{}}
\toprule\noalign{}
\begin{minipage}[b]{\linewidth}\raggedright
No.
\end{minipage} & \begin{minipage}[b]{\linewidth}\raggedright
what was wrong
\end{minipage} & \begin{minipage}[b]{\linewidth}\raggedright
how it was found
\end{minipage} & \begin{minipage}[b]{\linewidth}\raggedright
correction
\end{minipage} \\
\midrule\noalign{}
\endhead
\bottomrule\noalign{}
\endlastfoot
N-E1 & \textbf{Hard Rule 2 was broken.} From the author's instruction to
the first commit of the folder, about an hour and a half, nothing was
committed or carried to his folder. & The author pointed it out. & Every
step ends with a commit that is transferred to his folder at once
(Section 7, V9). \\
N-E2 & The first version of the search for the state of lowest energy
per particle refined only around the lowest value of a coarse scan of
the density. A narrow bound minimum between two scan points was missed,
and binding thresholds came out up to 2.7 per cent too high. & The
command \texttt{forms} prints beside every threshold the threshold of
the routine for matter of one kind; the two disagreed. & Commit
\texttt{6d654e8}: a finer scan, and a refinement around every local
minimum. \\
N-E3 & The folder began as a copy of another folder, so that its folder
\texttt{data} at first held the data files of the other program; a file
that no command of this folder wrote again would have carried the
numbers of the other program into the documents. & Foreseen when the
folder was created. & No file of the copy was committed before a command
of this folder had written it again. The script \texttt{check\_all.sh}
now verifies that every file in \texttt{data} names this folder in its
head. \\
N-E4 & Check \texttt{fm-F3} required a state at every polarization. In
form Q, at three and six times the threshold and at high density, 108 of
3936 states with unequal halves do not exist with non-negative masses. &
The first complete run of the command \texttt{forms} reported six of
nine checks. & The check states the fact and requires the states with
equal halves and with one filled half. \\
N-E5 & Check \texttt{fm-F4} covered clock readings up to 6. In form Q
the state with equal halves reaches the eigenvalue 1 of the matrix
\(\mathsf R_{\rm gap}\) at the clock reading 5.30 (\(D=3\), set A), and
states with slightly unequal masses appeared in the table. & The same
run. & The table ends at the clock reading 5 and carries the eigenvalue
as a column; the check reports where each case reaches 1. \\
N-E6 & Check \texttt{fm-F5} required that the pressure of a gas of fixed
density become negative in every case. In form C at the higher density
it stays positive in 3 of the 8 cases. & The same run. & The check
states what happens in each case. \\
N-E7 & The vacuum of the two halves has two solutions that differ by the
exchange of the halves. The first version of the command printed
whichever the solver found, and the first version of the independent
check compared the signed masses. & The independent check failed on a
case in which the two programs had found different members of the pair.
& The program prints a fixed representative (the half with the positive
mass first); the independent check compares the magnitudes and the
relative sign. \\
N-E8 & The gap equation was solved with the upper limit 60 of its
momentum integral, taken over from the earlier programs of the project;
the gaps depend on it. & The gap solver of the independent script used
another limit and disagreed with the program. & Commit \texttt{69d03e8}:
the ultraviolet criterion, the limit \(2\times10^4\), check
\texttt{pc-P7} (Section 9, F6 and F7). \\
N-E9 & \textbf{Hard Rule 5 was broken.} The first version of \texttt{P3}
(commit \texttt{9f8a206}) introduced symbols that had not been compared
with the MASTER registry; fifteen groups of them collided, or nearly
collided, with registered symbols. & Comparison of \texttt{P3} with the
registry before the next document was written. & Commit
\texttt{2ef3704}: the symbols entered in the registry, \texttt{P3}
rewritten, check MS8; \texttt{P1}, Section 3.3, lists what was
replaced. \\
N-E10 & A column of the data files had a name that is a symbol of one of
the author's notebooks (check MS3 failed). & The check. & The eight
names of the mean-field couplings in the program and in the data files
were changed; the four commands concerned were executed again; every
number is unchanged. \\
N-E11 & While \texttt{P3} was being rewritten: a script changed the
registry in memory and did not write the file; a replacement of the name
of a matrix throughout the text also changed a sentence about the letter
itself; a test for leftover symbols gave a false alarm on an index list.
& Reading the saved files and the differences. & Corrected before the
commit. \\
N-E12 & An equation of \texttt{P3} was wider than the page. A line of a
report of checks was too long for the typesetting program, which
stopped. The script \texttt{check\_all.sh} then compared the printed
reports with the unfolded files and called two documents outdated that
were not. & The builds; the script. & The equation is set on two lines;
the script that assembles the documents continues a report line of more
than 1800 characters on the next line; \texttt{check\_all.sh} folds in
the same way before it compares. \\
N-E13 & Drafts contained a number that no program states (in a reworded
sentence), a section number that the number check took for a result, the
name of a generated text that does not exist, and a reference to an
appendix by a number that had changed. & The number check; the script
that assembles the documents. & Corrected before the commits. \\
N-E14 & One transfer of a commit to the author's folder failed half-way:
the permission to replace files there had lapsed, and the program git
left a lock file and two files of the new commit as untracked files. &
The output of the transfer. & The permission was renewed by the author;
the lock was removed; the two files were compared by their hashes with
those of the commit, found identical, and removed; the transfer was
repeated. No file of the author's was touched. \\
N-E15 & A generated text counted the evolutions with gravity whose rates
run away over all sixteen evolutions of the folder, where the sentences
that use it speak of the eight of one spacetime. & Writing the section
on gravity of the manuscripts. & A count for each spacetime. \\
N-E16 & \textbf{A document with two outdated numbers was committed.}
Commit \texttt{25d5665} of the manuscripts quoted, in the section on the
effective action, two numbers of the comparison of the Ginzburg--Landau
coefficients that belong to the upper limit 60 (1.0 and 7.9 per cent;
the program of this folder states 1.2 and 9.8). The sentence had been
typed by hand in the predecessor of this folder. & The number check no
longer found the second number among the numbers that the programs
state. & Commit \texttt{a824eab}. Then every number that the programs of
the folder with the limit 60 state and those of this folder do not state
was searched for in the sources of all documents of this folder: no
other occurrence. \\
N-E17 & In the same commit the table of the checks in both manuscripts
gave 7 of 8 for the checks of the symbols, because \texttt{P1},
\texttt{P2} and a status file of the folder still contained replaced
symbols. The message of the commit says so. & Known when it was
committed. & The documents were rebuilt when \texttt{P1} and \texttt{P2}
were complete, and the status file was removed. \\
N-E18 & The 7+1 manuscript said that gaps computed with the limit 60 are
too small ``by the percentages'' of a table whose classes also contain
kernels of the sign that does not occur on the 7+1 spacetime. & Reading
the typeset page. & ``By up to''; the caption says that the classes hold
the kernels of both signs. \\
N-E19 & The first build of the manuscripts had seven boxes wider than
their space and one table of symbols taller than a page. & The build
script, which fails on them. & Equations and tables rearranged; the
table of symbols divided in two. \\
\end{longtable}

\textbf{The cause of N-E16, and the remedy.} The number check accepts a
number of the text if any program states that number anywhere; it cannot
know which statement a sentence means. A number typed by hand can
therefore survive a change of the calculation behind it. The remedy used
here is the comparison described in the row N-E16. Where a sentence
quotes a result of a calculation that may change, the manuscripts of
this folder take the number from a generated text and not from the
keyboard; the sentence of N-E16 was one of those that had been typed.

\hypertarget{errors-of-the-work-that-this-folder-inherited}{%
\subsection{12.2 Errors of the work that this folder
inherited}\label{errors-of-the-work-that-this-folder-inherited}}

These errors were made, and corrected, in the folder
\texttt{Revision\_Deflation\_CJT\_and\_Kohn-Sham}; the commits named are
commits of that work. For the errors C-E1 to C-E11, which were corrected
before anything was committed, the record kept during the work names the
error and what exposed it; an entry says no more than the record
contains.

\begin{longtable}[]{@{}
  >{\raggedright\arraybackslash}p{(\columnwidth - 6\tabcolsep) * \real{0.2500}}
  >{\raggedright\arraybackslash}p{(\columnwidth - 6\tabcolsep) * \real{0.2500}}
  >{\raggedright\arraybackslash}p{(\columnwidth - 6\tabcolsep) * \real{0.2500}}
  >{\raggedright\arraybackslash}p{(\columnwidth - 6\tabcolsep) * \real{0.2500}}@{}}
\toprule\noalign{}
\begin{minipage}[b]{\linewidth}\raggedright
No.
\end{minipage} & \begin{minipage}[b]{\linewidth}\raggedright
what was wrong
\end{minipage} & \begin{minipage}[b]{\linewidth}\raggedright
how it was found
\end{minipage} & \begin{minipage}[b]{\linewidth}\raggedright
correction
\end{minipage} \\
\midrule\noalign{}
\endhead
\bottomrule\noalign{}
\endlastfoot
C-E1 & The first formula for the pair couplings acted on pair states
without projecting on the states of positive energy and gave other
integers. & Comparison with the exact diagonalization for two particles
at rest. & The projected map; checks \texttt{71s-S9},
\texttt{44q-S9}. \\
C-E2 & An early text called the author's values \(G_S\), \(G_V\) of the
Kohn--Sham example the Hartree--Fock couplings of \(g_0\), \(g_1\) with
the same numbers. & The exact Hartree--Fock formula (check
\texttt{self-X1}). & The texts distinguish the two; finding F2. \\
C-E3 & In the first version of the check of the mass of the phase mode
(\texttt{uni-U3}) a deviation appeared that was an artefact of the
momentum cutoff. & The check failed. & Both sides of the comparison are
evaluated with the same cutoff. \\
C-E4 & In the history of deflating matter (\texttt{ks-K5}) a
minimization ran without bounds. & Unphysical values in the output. &
The search is bounded. \\
C-E5 & An angular average of the pairing amplitudes evaluated a function
outside its domain. & The program stopped. & Corrected. \\
C-E6 & The tolerance of an independent comparison of two integrations of
the deflating spacetime with matter (\texttt{Icks-G5}) was set tighter
than the accuracy of the comparison. & The check failed although both
integrations were correct. & The tolerance follows from the measured
accuracy, which the check states. \\
C-E7 & Labels of several new figures collided with curves, and some used
the names of the program or of the author's document instead of the
registered symbols. & Looking at each figure. & Figures redrawn. \\
C-E8 & The registry of symbols at first contained collisions introduced
by this sub-project: one letter for a velocity dispersion and for a
self-energy, one for a pair state and for a pair correlation, and two
definitions of \(x\). & The checks of the registry and reading the
manuscripts. & New symbols \(\sigma_v\), \(\Phi(\vec p)\); one
definition of \(x\). \\
C-E9 & A generated table contained section numbers typed by hand. &
Reading the built manuscript. & The table carries references that the
typesetting resolves. \\
C-E10 & The description of the test with deliberate errors repeated a
statement that belonged to a test of the first revision. & Comparison
with the report of the script of this folder. & The text states what the
report of this folder says. \\
C-E11 & A clause about the degeneracy of pair states in the section on
the Higgs hypothesis had not been verified. & Review of the section
against the checks. & Removed. \\
C-E12 & \textbf{A failing check was committed.} The message of commit
\texttt{a232561} says that the symbol checks pass \texttt{7\ of\ 7}; the
report in that commit says \texttt{6\ of\ 7}. With the documents
\texttt{P1} and \texttt{P2} in the folder, check MS4 failed: two
generated texts wrote the rates as plain text. & Reading the whole
report instead of its last lines, one step later. & Commit
\texttt{990d697}: both texts write the rates upright. \\
C-E13 & \textbf{A second failing check was committed.} Since \texttt{P2}
was added, the check of the numbers failed (it did not recognize
\texttt{Section\ 5.11} as a section number, and two run times had no
entry in its lists), and \texttt{P2} printed the failing report. & As
C-E12. & Commit \texttt{8231b87}. \\
C-E14 & \texttt{P2} printed a caption with an unresolved placeholder. &
The new script \texttt{check\_all.sh}. & Commit \texttt{8231b87}. \\
C-E15 & The running header of \texttt{P1} and \texttt{P2} named the
folder of the first revision. & Looking at the typeset pages. & Commit
\texttt{56f9162}. \\
C-E16 & In the figure of the creation of particles the labels of two
axes touched, and two legends lay on curves; the figure is in the 7+1
manuscript. & Looking at the typeset page of the new document. & Commit
\texttt{56f9162}: labels on two lines, legends moved, manuscript
rebuilt. \\
C-E18 & The headers of two generated tables carried equation numbers of
documents of the first revision. & Writing the appendix of complete
tables. & The headers describe the columns; the manuscripts refer to
their own equations. \\
C-E19 & In eleven of the thirty figures a label or a legend touched a
curve, a marked line or the frame (figures of the bound pair, of the
effective action, of the vacuum mass, of the extinction, of the
geodesics, of the two evolutions with gravity, of the Hubble constant,
of the quench, of the scalar modes, of the forecast, and of superfluid
dark matter); and three figures wrote the coordinates \(x0,\dots,x3\)
with two dots. Most of these figures had been taken over from the first
revision without being looked at again. & Every figure was looked at,
one by one, after C-E16. & Labels and legends moved, lines shortened,
captions adjusted where a label was shortened; all documents rebuilt. \\
C-E17 & A sentence drafted for ``dirac matrices deflating'' took the
roots of the constraint of the gravitational equations for rates of
growth. & Reading the statement of check \texttt{gr-G5} before the
document was committed. & The document says what the check says. \\
\end{longtable}

\textbf{The cause of C-E12 to C-E14, and the remedy.} After long builds
only the last lines of the outputs were read, and a verdict further up
was missed. The script \texttt{code/scripts/check\_all.sh} now runs the
checks of the documents, reads the verdict of every report in
\texttt{data}, verifies that no built document contains an unresolved
placeholder and that every report printed in a document is the current
one, and prints \texttt{ALL\ PASS} only then. Nothing is committed
without it.

\hypertarget{the-symbols-used-in-this-document}{%
\section{13. The symbols used in this
document}\label{the-symbols-used-in-this-document}}

\begin{longtable}[]{@{}
  >{\raggedright\arraybackslash}p{(\columnwidth - 2\tabcolsep) * \real{0.5000}}
  >{\raggedright\arraybackslash}p{(\columnwidth - 2\tabcolsep) * \real{0.5000}}@{}}
\toprule\noalign{}
\begin{minipage}[b]{\linewidth}\raggedright
symbol
\end{minipage} & \begin{minipage}[b]{\linewidth}\raggedright
meaning
\end{minipage} \\
\midrule\noalign{}
\endhead
\bottomrule\noalign{}
\endlastfoot
\(x0,\dots,x7\); \(xt\) & the eight coordinates (\(x4\): ordinary time);
a deflating coordinate, \(t=5,6,7\) \\
\(\mathrm{g55}\), \(\mathrm{g66}\), \(\mathrm{g77}\);
\(\mathrm{rate}_t\); \(\mathrm{gsum}\) & the three rates of deflation
(the author's names; upright, never with a subscript); the rate of
\(xt\); their sum \\
\(\mathrm{const}_t\) & the positive constants of the hypothesis, one for
each deflating coordinate \\
\(g^{(4+4)}_{\mu\nu}\), \(g^{(7+1)}_{\mu\nu}\) & the metrics of the two
deflating spacetimes, space-like components positive \\
\(V_d\) & \(\sqrt{\lvert\det g\rvert}\), the proper volume of a unit
coordinate cell of \(x5,x6,x7\); proportional to
\(e^{-\mathrm{gsum}\,x4}\) \\
\(\tau_d\) & the deflation clock, \(\mathrm{gsum}\,(x4-x4_{\rm ref})\),
with a reference time \(x4_{\rm ref}\) \\
\(\psi\); \(\hat\psi\) & the sixteen-component Dirac field; the weighted
field \(\sqrt{V_d}\,\psi\) \\
\(\gamma^A\); \(\gamma^\mu\); \(\nabla_\mu\) & the eight real
\(16\times16\) Dirac matrices of an orthonormal frame, \(A=0,\dots,7\);
the same with a coordinate index; the covariant derivative of a
spinor \\
\(\beta\); \(\bar\psi\); \(\alpha^j\); \(\mathsf h\) & \(\gamma^4\) on
7+1 and \(-i\gamma^4\) on 4+4; the adjoint \(\psi^\dagger\beta\); on 4+4
the matrices \(-\gamma^4\gamma^j\), \(j=0,1,2,3\); the Hamiltonian
matrix of a free particle \\
\(C^-\); \(\tilde\psi\); \(W_{1644}\) & on 4+4: the matrix
\(\gamma^0\gamma^1\gamma^2\gamma^3\); the covariant adjoint
\(\psi^\dagger C^-\); the matrix \(iC^-\gamma^4\), Hermitian with the
eigenvalues \(+1\) and \(-1\), eight times each; its two eigenspaces are
the two \emph{halves} of the field \\
\(\mathrm{Adj}\) & the matrix of the adjoint in a bilinear
\(\psi^\dagger\mathrm{Adj}\,\Theta^K\psi\): \(\beta\) on 7+1 and in form
Q, \(C^-\) in form C \\
form C, form Q & the two forms of the self-interaction of the 4+4 model
in the quantized sector: with the covariant adjoint (the definition) and
with the adjoint of the positive Lagrangian (the alternative) \\
\(\gamma_{44q}^A\); \(\Theta_Q^K\); \(\mathrm{sgn}_K\),
\(\mathrm{sgn}_K^Q\) & the Dirac matrices of form Q, \(-i\gamma^A\) for
\(A\le4\) and \(\gamma^A\) for \(A\ge5\); their products with a phase
that makes the bilinears real; the signs \(\pm1\) of the squares of
\(\Theta^K\) and of \(\Theta_Q^K\) \\
\(I_r^{\rm C}\), \(I_r^{\rm Q}\); \(\mathcal L_{\rm self}^{\rm C}\),
\(\mathcal L_{\rm self}^{\rm Q}\) & the nine invariants and the
self-interaction in the two forms \\
\(g_0^{\rm eff}\), \(g_1^{\rm eff}\), \(g_7^{\rm eff}\),
\(g_8^{\rm eff}\); \(g\); set A, set B & the four independent
combinations of the nine couplings (scalar, vector, axial vector,
pseudoscalar); the strength of the couplings; the two sets of couplings
\(g_0=-g\), and \(g_0=-g\) with \(g_1=+0.1\,g\) \\
\(\mathsf R_{\rm gap}\) & for matter of the two halves, the matrix of
the derivatives of the right-hand sides of the two gap equations with
respect to the two masses \\
\(s_{\rm ct}\); \(m_{\rm ct}\); \(g_{VS}\) & the sign \(\pm1\) with
which the contact couplings enter the pairing potential of a pair; their
range parameter; the coupling of the Vorton-Scalar to the fermions \\
\(U_{\rm uv}\); \(u_S\), \(u_V\) & the ultraviolet strength of the
kernel of the gap equation; its two numerical constants \\
\(D\); \(\tau_d\) in the tables of Section 6 & the number of space
dimensions in which the particles have momentum (4, or 3 for fields in
ordinary space); the deflation clock \\
\(m\); \(m_D\) & the mass of the fermion; the part of it that is a bare
mass in the model with the Vorton-Scalar \\
\(\mathcal L_{\rm kin}\), \(\mathcal L_{\rm mass}\),
\(\mathcal L_{\rm self}\) & the three parts of the Lagrangian: kinetic
term, mass term, self-interaction \\
\(K\); \(r\); \(\Theta^K\), \(\Theta_K\) & a strictly increasing list of
indices out of \(0,\dots,7\); its length, the rank; the product of the
Dirac matrices of the list, with a phase that makes the bilinears real,
and its inverse \\
\(I_r\); \(g_r\) & the invariant of rank \(r\) of the covariant
Lagrangian; the nine couplings, \(r=0,\dots,8\) (\(g_0\) scalar, \(g_1\)
vector) \\
\(\mathsf F\) & the Fierz matrix, a \(9\times9\) matrix of rational
numbers with \(\mathsf F^2=1\) \\
\(G_S\), \(G_V\); \(G_S^{\rm HF}\), \(G_V^{\rm HF}\) & the scalar and
the vector coupling of the mean-field energy of uniform matter; their
values from the nine couplings in the Hartree--Fock approximation \\
\(k_F\) & the Fermi momentum \\
\(\Lambda\) & a momentum cutoff \\
\(SO(4)\); \(U(1)_E\) & for fields that depend on \(x1,x2,x3,x4\) only:
the rotations of the four extra axes \(x0,x5,x6,x7\), generated by the
products \(\gamma^a\gamma^b\) of their Dirac matrices; a second phase
symmetry, generated by a multiple of
\(\gamma^0\gamma^5\gamma^6\gamma^7\) \\
Vorton-Scalar & the complex scalar field of the model, coupled to
\(\bar\psi\psi\); its radial mode and its phase mode are its two kinds
of excitation \\
\end{longtable}

\end{document}
